\definition \textbf{Zufallsvariable}
$$
    X: \Omega \to \R \qquad \text{s.d} \quad \forall a \in \R:\quad \Bigl\{ \omega \in \Omega\ \Big|\ X(\omega)\leq a \Bigr\} \in \F
$$
\subtext{$(\Omega, \F, \P)$ ein Wahrscheinlichkeitsraum}

\footnotesize
\definition \textbf{Indikatorfunktion}
$$
    \forall \omega \in \Omega:\quad \mathbb{I}_A(\omega) := \begin{cases}
        0 & \omega \notin A \\
        1 & \omega \in A
    \end{cases}
    \color{gray}
    \qquad A \in \F
    \color{black}
$$
\normalsize

\notation Ereignisse bezüglich Zufallsvariablen
\begin{align*}
    & \{ X \leq a \}      &=\quad& \Bigl\{ \omega \in \Omega \ \Big|\ X(\omega) \leq a \Bigr\} \\
    & \{ a < X \leq b \}  &=\quad& \Bigl\{ \omega \in \Omega \ \Big|\ a < X(\omega) \leq b \Bigr\} \\
    & \P[X \leq a]        &=\quad& \P\Bigl[\{ X \leq a \}\Bigr]
\end{align*}

\definition \textbf{Verteilungsfunktion} $F_X: \R \to [0,1]$
$$
    \forall a \in \R:\qquad F_X(a) = \P[X \leq a]
$$

\theorem \textbf{Eigenschaftern der Verteilungsfunktion}

\begin{tabular}{ll}
    (i)     & $F_X$ monoton \\
    (ii)    & $F_X$ rechtsstetig \\
    (iii)   & $\underset{a \to -\infty}{\lim} F_X(a) = 0 \quad\land\quad \underset{a \to \infty}{\lim}F_X(a)=1$
\end{tabular}

\definition \textbf{Unabhängigkeit}\\
\smalltext{$X_1,\cdots,X_n \text{ unabhängig } \iffdef§  \forall x_1,\cdots,x_n \in \R:$}
$$
    \P[X_1 \leq x_q,\cdots, X_n \leq x_n] = \P[X_1 \leq x_1] \cdots \P[X_n \leq x_n]
$$

\theorem \textbf{Unabhängigkeit von Gruppierungen}\\
\smalltext{$X_1,\cdots X_n$ sind unabhängig, dann sind auch $Y_1,\cdots,Y_k$ unabhängig:}
$$
    Y_1 = \phi_1(X_1,\cdots,X_{i_1}), \cdots , Y_k = \phi_k(X_{i_{k-1}+1},\cdots,X_{i_k})
$$
\subtext{$1 \leq i_1 < i_2 < \cdots < i_k \leq n \text{ sind Indizes},\quad \phi_1,\cdots,\phi_k \text{ sind Abbildungen}$}

\newpage

\definition \textbf{Folgen von Zufallsvariablen}

\begin{tabular}{llll}
    (i)     & unabhängig & $\iffdef$ & $\forall n \geq 1: X_1,\cdots,X_n \text{ unabhängig}$ \\
    (ii)    & i.i.d      & $\iffdef$ & unabhängig, und $\forall i, j: F_{X_i} = F_{X_j}$
\end{tabular}

\subtext{i.i.d = Independent \& identically distributed}

\subsection{Diskrete Zufallsvariablen}

\definition \textbf{Diskrete Zufallsvariable}
$$
    X: \Omega \to \R \text{ diskret } \iffdef \exists (W \subset \R) \preceq \N:\quad \P[X \in W] = 1
$$
\subtext{bzw. die Werte von $X$ liegen f.s. in $W$}

\lemma \textbf{Variablen diskreter Grundräume sind diskret}
$$
    \Omega \preceq \N \implies X: \Omega \to \R \text{ ist diskret}
$$

\definition \textbf{Verteilung diskreter Variablen}
$$
    \Bigl( p(x) \Bigr)_{x \in W} \text{ s.d. } p(x) := \P[X = x] \text{ heisst Verteilung}
$$
\subtext{$X$ ist diskret mit $W \preceq \N\quad p(x)$ ist \textit{nicht} $F_X$}

\theorem $\forall p(x):\quad \displaystyle\sum_{x \in W} p(x) = 1$ \subtext{$\quad p(x)$ ist eine diskrete Verteilung}

\lemma \textbf{Zur diskreten Verteilung existiert eine Variable}
$$
    \forall \Bigl( p(x) \Bigr)_{x \in W} \in [0,1]:\quad \exists (\Omega, \F, \P), X \text{ mit Verteilung } p(x)
$$
\subtext{D.h man kann sagen: "Sei $X$ eine Variable mit Verteilung $\bigl(p(x)\bigr)_{x \in W}$"}

\theorem \textbf{Diskrete Verteilungsfunktion}
$$
    F_X(x) = \P[x \leq X] = \sum_{y \leq x} p(y) \qquad y \in W
$$

\newpage

\subsection{Diskrete Verteilungen}

\definition \textbf{Bernoulli-Verteilung} $X \sim \text{Ber}(p)$\\
\smalltext{Intuitiv: Münzwurf}
$$
    \P[X=1] = p \qquad \P[X=0] = 1-p
$$
\subtext{$0 \leq p \leq 1,\quad W = \{0,1\}$}

\definition \textbf{Binomial-Verteilung} $X \sim \text{Bin}(n,p)$\\
\smalltext{Intuitiv: Anzahl Erfolge bei wiederholtem Bernoulli-Experiment}
$$
    \P[X=k] = \binom{n}{k}p^k(1-p)^{n-k}
$$
\subtext{$0 \leq p \leq 1,\quad n \in \N,\quad W = \{0,\ldots, n\}$}

\theorem \textbf{Bernoulli-Summen sind Binomial-verteilt}
$$
    S_n := X_1 + \ldots + X_n \sim \text{Bin}(p, n)
$$
\subtext{$0\leq p\leq n,\quad n \in N,\quad X_1,\ldots,X_n \sim \text{Ber}(p) \text{ unabhängig}$}

\scriptsize
\lemma $\text{Bin}(1,p) = \text{Ber}(p)$

\lemma $X_1 \sim \text{Bin}(n, p), X_2 \sim \text{Bin}(m, p) \implies X_1 + X_2 \sim \text{Bin}(n+m,p)$

\lemma \textbf{Binomialverteilung erfüllt die Summenvoraussetzung}
$$
    \sum_{k=0}^{n} p(k) = \sum_{k=0}^{n} \binom{n}{k} p^k(n-p)^{n-k} = (p + 1 - p)^n = 1
$$
\normalsize

\definition \textbf{Geometrische Verteilung} $X \sim \text{Geom}(p)$\\
\smalltext{Intuitiv: Bernoulli Experiment erfolgreich beim $k$-ten Versuch}
$$
    \P[X=k] = p\cdot(1-p)^{k-1}
$$
\subtext{$0 < p \leq 1,\quad k \in \N_{\neq0}$}